\chapter{Discussione: Sintesi Trasversale, Confronto Tecnologico e Limiti}
\label{chap:discussion}

% ============================================================================
\section{Sintesi Trasversale tra Sorveglianza UE, Richiami USA e Monitoraggio Chimico}
\label{sec:disc_synthesis}
% ============================================================================

I tre moduli di DashMe analizzano fonti con giurisdizioni, unità di analisi e periodi diversi: RASFF per le notifiche europee, FDA per i richiami statunitensi ed EFSA / WHO GEMS per i controlli chimici \citep{rasff_portal,fda_recalls_portal,who_gems_efsa}. Poiché i record non sono collegati tra loro, il confronto individua analogie e differenze tra prospettive, ma non ricostruisce un unico ciclo causale del rischio.

Una prima convergenza riguarda la distinzione tra frequenza e gravità. Nei dati FDA i settori con più prodotti richiamati non coincidono con quelli con maggiore quota di Classe~I; in RASFF, i volumi di notifica non descrivono da soli la quota di rischio grave; nei dati di laboratorio, il tasso di positività dipende anche dalle strategie di campionamento. Questi indicatori rispondono a domande diverse e non sono direttamente confrontabili come misure di prevalenza.

Le analisi RASFF e FDA segnalano entrambe la rilevanza dei pericoli microbiologici, ma in contesti distinti: \textit{Salmonella spp.} è il contaminante più frequente nel dataset RASFF, mentre nel campione FDA \textit{Listeria} emerge soprattutto tra i richiami associati ai prodotti refrigerati e surgelati. L'associazione osservata tra cause, temperatura e packaging è descrittiva e non dimostra, da sola, un nesso causale.

Le viste geografiche e operative mettono in luce criticità diverse: RASFF evidenzia flussi di notifica e respingimenti frontalieri, mentre FDA documenta richiami ancora aperti e reti distributive di diversa ampiezza. Questi risultati suggeriscono di affiancare al monitoraggio delle allerte indicatori di completamento dei richiami, senza assumere che le notifiche europee e statunitensi riguardino gli stessi prodotti o eventi.

Infine, EFSA / WHO GEMS fornisce una prospettiva preventiva sui controlli di laboratorio. Le positività al cadmio nelle frattaglie e l'elevato numero di analisi generato dai pannelli multi-residuo mostrano come la selezione delle matrici e dei protocolli possa orientare il monitoraggio; non permettono però di dedurre un collegamento diretto con specifiche allerte RASFF o richiami FDA.

% ============================================================================
\section{Confronto Critico tra Qlik Sense, Tableau e Power BI}
\label{sec:disc_platforms}
% ============================================================================

L'implementazione parallela del progetto DashMe su tre piattaforme leader di mercato consente di valutare in modo empirico i compromessi architetturali tra Qlik Sense, Tableau e Microsoft Power BI rispetto ai principi di progettazione dimensionale e visiva~\citep{kimball2013,few2013,sharda2020}. La Tabella~\ref{tab:platform_comparison} sintetizza il confronto lungo cinque dimensioni tecniche e operative.

\begin{table}[htbp]
\centering
\caption{Confronto tecnico-funzionale tra le tre piattaforme di Business Intelligence impiegate nel progetto DashMe.}
\label{tab:platform_comparison}
{\small
\begin{tabularx}{\textwidth}{@{}>{\raggedright\arraybackslash}p{2.5cm}>{\raggedright\arraybackslash}X>{\raggedright\arraybackslash}X>{\raggedright\arraybackslash}X@{}}
\toprule
\textbf{Dimensione} & \textbf{Qlik Sense (RASFF)} & \textbf{Tableau (FDA)} & \textbf{Power BI (EFSA/WHO)} \\
\midrule
\textbf{Motore dati e architettura} &
Motore associativo \textit{in-memory} (QIX) basato su indicizzazione relazionale e stato selezioni verde-bianco-grigio~\citep{qlik_docs}. &
Motore colonnare \textit{Hyper} (estrazione \texttt{.hyper}) o connessione \textit{Live}, ottimizzato per query visive interattive~\citep{tableau_docs}. &
Motore colonnare \textit{VertiPaq} su modello \textit{Star Schema} alimentato da ETL in \textit{Power Query M}~\citep{powerbi_docs}. \\
\addlinespace[0.35em]
\textbf{Linguaggio analitico e ETL} &
Script di caricamento (\texttt{WildMatch}, pulizia inline) ed espressioni con \textit{Set Analysis} per metriche contestuali. &
Campi calcolati condizionali, funzioni stringa (\texttt{CONTAINS}) e aggregazioni distinte (\texttt{CNTD}) a livello di vista. &
Trasformazioni strutturate in M e misure riusabili in \textit{DAX} (\texttt{CALCULATE}, gestione nulli e fasce). \\
\addlinespace[0.35em]
\textbf{Espressività visiva e UX} &
Forte esplorazione multidimensionale su quattro fogli coordinati; treemap, matrici a bolle, mappe e diagramma di Sankey. &
Massima flessibilità compositiva e narrativa: mappe, barre impilate, \textit{packed bubbles} e matrici di densità multi-livello. &
Layout rigido ma coerente su griglia; \textit{Azure Maps}, \textit{Ribbon}, \textit{Waterfall}, \textit{Gauge} e \textit{Decomposition Tree}. \\
\addlinespace[0.35em]
\textbf{AI integrata e analytics avanzata} &
Assente nelle viste standard sviluppate; analisi guidata dalla scoperta associativa dei valori esclusi. &
Funzionalità statistiche opzionali (trend, cluster) non richieste dal flusso causale basato su regole esplicite. &
Componenti nativi di ML (\textit{Key Influencers} e \textit{Segmenti Principali} ad albero decisionale) integrati nel report. \\
\addlinespace[0.35em]
\textbf{Trade-off operativi} &
Ottimo per individuare correlazioni impreviste e flussi di rete; più sensibile a sporcizia CSV in fase di ingestione. &
Eccellente prototipazione visiva rapida; logica \textit{rule-based} nei campi calcolati soggetta a collisioni lessicali se non testata. &
Elevata scalabilità su $>2,4\text{Mln}$ di righe e governance rigorosa delle misure; richiede disciplina di modellazione dimensionale. \\
\bottomrule
\end{tabularx}
}
\end{table}

Dal confronto emergono implicazioni progettuali chiare per la scelta dello strumento in contesti di \textit{Risk Management} e sorveglianza regolatoria:
\begin{itemize}
    \item \textbf{Esplorazione associativa contro interrogazione gerarchica}: il motore QIX di Qlik Sense si rivela particolarmente efficace nell'analisi di tracciabilità e geopolitica RASFF perché preserva visivamente tanto i valori associati quanto quelli esclusi da una selezione, facilitando la scoperta rapida di corridoi anomali (come l'auto-notifica intra-UE rispetto ai gateway extra-UE). Per contro, l'assenza di una fase di tipizzazione e validazione rigida analoga a Power Query rende lo script di caricamento più vulnerabile a irregolarità sintattiche dei file testuali grezzi.
    \item \textbf{Espressività visiva contro manutenibilità semantica}: Tableau offre la maggiore agilità nella composizione di dashboard dense e ad alto impatto comunicativo (come la sequenza di tre viste complementari sui richiami FDA, capace di bilanciare conteggi distinti \texttt{CNTD(ID Evento)} sulle mappe e volumi di prodotto sulle cause). Tuttavia, la comodità di implementare la tassonomia direttamente nei campi calcolati del workbook espone il modello a fragilità di manutenzione --- quali la collisione cromatica su \texttt{\#529ee0} o i conflitti tra sottostringhe --- quando le regole non vengono incapsulate in una pipeline ETL separata.
    \item \textbf{Scalabilità colonnare, DAX e AI descrittiva}: Power BI dimostra la superiore tenuta architetturale su volumi superiori ai due milioni di righe, grazie alla separazione netta tra preparazione dei dati in Power Query M, modellazione a stella in VertiPaq~\citep{kimball2013} e definizione centralizzata delle misure DAX (fondamentale per escludere i 3.330 valori nulli dal denominatore del tasso di positività). L'integrazione nativa di \textit{Decomposition Tree} e \textit{Key Influencers} riduce inoltre la distanza tra Business Intelligence descrittiva e analisi diagnostica automatizzata~\citep{sharda2020}, permettendo al decisore non solo di osservare la positività media, ma di isolarne istantaneamente i moltiplicatori statistici (2,53x) e i segmenti ad alta incidenza (fino al 72,9\%).
\end{itemize}

% ============================================================================
\section{Implicazioni per il Project Management e il Risk Management}
\label{sec:disc_implications}
% ============================================================================

Le evidenze suggeriscono quattro priorità operative per i responsabili di progetto, qualità e \textit{supply chain}:

\begin{enumerate}
    \item \textbf{Controlli frontalieri basati sul rischio}: concentrare le verifiche documentali e microbiologiche sui flussi e sulle combinazioni matrice-pericolo che le analisi RASFF identificano come più critiche, affiancando ai controlli uniformi criteri di priorità aggiornabili.

    \item \textbf{Governance della chiusura dei richiami}: gestire ogni campagna come un progetto con responsabili, scadenze, conferme di ricezione dai distributori e riconciliazione delle scorte ritirate o distrutte; prevedere soglie di escalation per i richiami ancora aperti e per le reti multi-giurisdizionali.

    \item \textbf{Presidio HACCP per temperatura e confezionamento}: usare i risultati come ipotesi per orientare verifiche mirate su catena del freddo, sanificazione, integrità delle confezioni ed etichettatura. Le associazioni osservate sono descrittive e vanno verificate nel contesto operativo prima di modificare i piani di controllo.

    \item \textbf{Piani di campionamento laboratoristico}: valutare l'impiego di pannelli multi-residuo e dare priorità alle matrici che le analisi indicano come rilevanti, verificando la comparabilità metrologica prima di confrontare concentrazioni o definire soglie.
\end{enumerate}

% ============================================================================
\section{Limiti Metodologici e Qualità dei Dati}
\label{sec:disc_limits}
% ============================================================================

Un'interpretazione rigorosa dei risultati richiede di esplicitare le criticità di qualità del dato, di modellazione e di resa visiva emerse durante la validazione tecnica delle tre suite analitiche:

\begin{itemize}
    \item \textbf{Limiti del modulo Qlik Sense / RASFF}:
    \begin{enumerate}
        \item \textit{Fragilità di parsing CSV}: l'ingestione del file \texttt{RASFF\_window.csv} ha evidenziato 4 righe malformate contenenti doppi apici non gestiti nel campo testuale \texttt{subject}, determinandone l'esclusione in fase di caricamento (11.787 notifiche validate su 11.788 record grezzi).
        \item \textit{Residui di mancata classificazione}: nella treemap dei Top 15 contaminanti specifici della Dashboard~3 permangono 139 occorrenze etichettate con il segnaposto \texttt{-} (dato non classificato o assente al 14\textdegree{} posto).
        \item \textit{Asimmetria tra legenda e serie visualizzate}: il grafico multi-linea sulla stagionalità mensile espone in legenda 8 famiglie di pericolo a fronte delle sole 5 famiglie prevalenti effettivamente tracciate in ciascun mese.
        \item \textit{Concentrazione diacronica}: oltre il 98,8\% delle notifiche valide si concentra nella finestra 2020--2025, mentre la scarsa densità campionaria del periodo 2012--2019 produce oscillazioni estreme (0\%--100\%) della \textit{\% Rischio Grave} prive di stabilità statistica.
    \end{enumerate}

    \item \textbf{Limiti del modulo Tableau / FDA}:
    \begin{enumerate}
        \item \textit{Ampiezza delle classi residuali da parsing testuale}: la classificazione tramite espressioni condizionali su testo libero genera categorie residuali molto ampie, pari al 58,05\% per \texttt{Altro} in \texttt{Categoria Ragione di Richiamo} (60.681 prodotti su 104.537), al 72,26\% per \texttt{Altro / Non specificato} in \texttt{Tipo Packaging} e al 92,60\% per \texttt{Temperatura Ambiente / Altro} in \texttt{Modalità Conservazione}.
        \item \textit{Collisione lessicale ``peanut''/``nut'' negli allergeni}: nella logica di conteggio di \texttt{Dettaglio Allergene} e \texttt{Numero Allergeni Coinvolti}, la verifica simultanea delle sottostringhe \texttt{"peanut"} e \texttt{"nut"} fa sì che ogni menzione di arachide attivi entrambi i flag, dirottando 1.046 richiami mono-allergene da arachidi (su 1.600 menzioni complessive) nella bolla \textit{Più Allergeni Contemporaneamente} (2.766) ed eliminando visivamente la bolla autonoma delle arachidi.
        \item \textit{Collisione cromatica nella Dashboard~3}: la condivisione del medesimo codice colore azzurro (\texttt{\#529ee0}) tra \textit{adulterazioni farmaci} e \textit{Salmonella} nella vista dedicata al packaging può indurre letture visive fuorvianti della barra \textit{Vetro / Bottiglia / Barattolo} (in realtà dominata al 74,4\% da adulterazioni farmacologiche).
        \item \textit{Residua frammentazione anagrafica multi-sito}: nonostante la normalizzazione operata in \texttt{Ragione Sociale (Pulita)}, l'operatore \textit{Medline Industries LP} rimane suddiviso nella Top 15 tra \textit{Medline Industries LP - Northfield} (822 prodotti) e \textit{Medline Industries LP} (557 prodotti), anziché risultare consolidato in un unico gruppo da 1.379 prodotti.
    \end{enumerate}

    \item \textbf{Limiti del modulo Power BI / EFSA-WHO}:
    \begin{enumerate}
        \item \textit{Eterogeneità metrologica (\texttt{UnitID = 1}) sul P95}: nel calcolo del 95\textdegree{} percentile delle concentrazioni positive in Dashboard~1, l'inclusione di record con \texttt{UnitID = 1} non convertiti in $\mu\text{g/kg}$ distorce verso l'alto la dimensione della bolla riferita alla \textit{Francia}.
        \item \textit{Impatto della classe \texttt{Non Specificato (Other)} e dei valori nulli}: il dataset contiene 702.718 righe prive di analita specifico (escluse dal filtro di pagina nella sola Dashboard~2) e 3.330 record con \texttt{ResultValue} nullo (esclusi dal denominatore delle misure di positività), determinando basi statistiche differenti tra le pagine (39,63\% globale, 43,44\% nazionale e 36,96\% sul sottoinsieme identificato).
        \item \textit{Bias di campionamento mirato}: i tassi di positività elevati osservati in alcuni Paesi --- quali i Paesi nordici (68\%--85\%), l'\textit{Ucraina} prossima al 100\% su volumi ridotti o la \textit{Svizzera} identificata con moltiplicatore 2,53x nei \textit{Key Influencers} --- riflettono piani di sorveglianza mirati su matrici critiche e non vanno interpretati come stime di prevalenza casuale sull'intero mercato alimentare.
        \item \textit{Disallineamento di versioning tra esportazioni e modello}: l'introduzione del campo \texttt{Gruppo Alimentare (breve)} nel file \texttt{.pbix} al commit \texttt{c693e38} comporta un lieve disallineamento nominale rispetto alle etichette estese presenti nel precedente artefatto esportato (\texttt{docs/\-PowerBI/\-dash1\_2.pdf}).
    \end{enumerate}
\end{itemize}
